\documentclass[11pt]{article}
\usepackage[margin=1in]{geometry}
\usepackage{hyperref}
\usepackage{enumitem}
\title{LEGO Game and Conflict Resolution}
\author{Piter Garcia}
\date{\today}
\begin{document}
\maketitle

\section*{Grade and Class Match}
\textbf{Grade:} Grade 1\\
\textbf{Likely blocks:} Day 3 · 12:25-1:15 · Gargana; Day 3 · 2:15-3:05 · Polfleit; Day 5 · 12:25-1:15 · Montgomery

\section*{Lesson Focus}
Partner play, turn-taking, and problem solving during a hands-on build game.

\section*{Learning Targets}
\begin{itemize}[leftmargin=*]
\item Students follow basic partner and turn-taking expectations.
\item Students participate in a shared hands-on challenge.
\item Students reflect on how their group solved a conflict or problem.
\end{itemize}

\section*{Materials}
\begin{itemize}[leftmargin=*]
\item LEGO or similar build materials
\item Visual directions
\item Reflection prompt
\end{itemize}

\section*{Suggested Lesson Flow}
\begin{enumerate}[leftmargin=*]
\item Open with the exact device, tool, or routine students need in order to start successfully.
\item Model one example task before students begin independent or partner work.
\item Give students time to build, code, design, or document their work while circulating for support.
\item Close with a short share-out, reflection, or debugging discussion.
\end{enumerate}

\section*{Source Notes}
This lesson packet is enriched with transcript-derived lesson notes, the school schedule roster, and official starting-point resources. See the paired Markdown guide for clickable links.

\bibliographystyle{plain}
\bibliography{lego-game-and-conflict-resolution}
\end{document}
